\subsection{Прагмы GCC: битсеты и циклы}
\begingroup
\footnotesize
\let\latintexttt\texttt
\renewcommand{\texttt}[1]{{\fontencoding{T1}\selectfont\latintexttt{#1}}}
\begin{otherlanguage}{russian}
Для GCC на x86 с поддержкой AVX2 и POPCNT можно проверить
такой набор \textbf{в самом начале файла, до include}:
\begin{lstlisting}
#pragma GCC optimize("O3,unroll-loops")
#pragma GCC target("avx2,popcnt")
#include <bits/stdc++.h>
\end{lstlisting}
Это четыре опции в двух строках, а не универсальное ускорение.
\begin{itemize}[nosep,leftmargin=*]
\item \texttt{popcnt}: аппаратный подсчёт единиц, полезен для
\texttt{bitset::count()} и \texttt{\_\_builtin\_popcountll}.
\item \texttt{avx2}: разрешает 256-битные SIMD-инструкции.
Может ускорить массовые AND/OR/XOR; сам по себе не добавляет
векторную инструкцию popcount.
\item \texttt{O3}: более агрессивная оптимизация циклов и
векторизация. Современный GCC умеет векторизовать и при O2.
\item \texttt{unroll-loops}: разворачивание циклов; уменьшает
накладные расходы, но увеличивает код и иногда замедляет.
\end{itemize}
Начать можно с одного \texttt{target("popcnt")}, затем сравнить
полный набор на максимальном тесте.
\texttt{bmi,bmi2,lzcnt} добавляют отдельные битовые инструкции;
для обычного bitset не обязательны.
\texttt{fma} относится к совмещённому умножению и сложению, не к битсетам.

\textbf{Ограничения.}
Неподдерживаемый target может дать Illegal instruction на судье;
наличие инструкции на своём ПК этого не исключает.
Не включать \texttt{Ofast}/\texttt{fast-math} автоматически:
они ослабляют правила вещественной арифметики, включая
NaN/Inf и порядок вычислений. Для геометрии особенно опасно.
\texttt{O3} не включает \texttt{fast-math}.
Проверять предупреждения \texttt{-Wall -Wextra -Wpragmas}.

Для отдельного участка:
\begin{lstlisting}
#pragma GCC push_options
#pragma GCC optimize("O3")
#pragma GCC target("popcnt")
// определения выбранных функций
#pragma GCC pop_options
\end{lstlisting}
Для одного цикла: \texttt{\#pragma GCC unroll 4} перед ним.
\texttt{\#pragma GCC ivdep} обещает отсутствие зависимостей,
мешающих SIMD; нельзя ставить на произвольный DP-цикл.

Не переносить старые хаки, переопределяющие
\texttt{\_\_cplusplus} или \texttt{unsigned}, ради 64-битного
bitset. Размер блока зависит от ABI и библиотеки. При необходимости
контролировать его явно через \texttt{vector<uint64\_t>}.
Для битсетов \texttt{\& | \^{} << >>} работают блоками, а
\texttt{operator[]} в цикле перебирает отдельные биты.
{\scriptsize Источники:
\href{https://codeforces.com/blog/entry/129283}{CF: прагмы GCC},
\href{https://codeforces.com/blog/entry/77480}{CF: bitset},
\href{https://gcc.gnu.org/onlinedocs/gcc/Function-Specific-Option-Pragmas.html}{документация GCC}.
Подробнее: \texttt{doc/pragmas.md}.}

\end{otherlanguage}
\endgroup
